% ! TEX root = ../ag-19-20.tex

\chapter{Parțiale 2018--2019}

\textbf{Seria A -- Numărul 1}

1. (a) Fie subspațiile vectoriale:
\begin{align*}
  V_1 &= \{ (x, y, z) \in \RR^3 \mid 2x - y + 3z = 0 \} \\
  v_2 &= \dr{Sp}\{ (1, -1, 3), (1, 2, 0) \}
\end{align*}
ale spațiului vectorial $ \RR^3 $.

Să se determine cîte o bază și dimensiunea pentru $ \RR $-spațiile vectoriale
$ V_1, V_2, V_1 + V_2, V_1 \cap V_2 $. Să se determine coordonatele vectorului
$ v = (5, 9, 13) $ după o bază a lui $ V_1 + V_2 $.

(b) Fie $ U, V \seq \RR^{20} $. Să se arate că $ V \simeq U \Leftrightarrow \dim U = \dim V $.

\vspace{0.5cm}

2. (a) Fie $ f : \RR^3 \to M_2(\RR) $ o aplicație $ \RR $-liniară, definită prin:
\[
  f(a, b, c) = %
  \begin{pmatrix}
    a + 2c & 0 \\
    a - b & b + 2c
  \end{pmatrix}.
\]

Să se determine cîte o bază și dimensiunea pentru $ \dr{Ker} f $ și $ \dr{Im}f $, matricea
asociată lui $ f $ în bazele canonice și să se determine $ V \seq \RR^3 $ astfel
încît $ \dr{Ker}f \oplus V \simeq \RR^3 $.

(b) Se consideră $ \RR^4 $ un spațiu vectorial și o aplicație liniară
$ f : \RR^4 \to \RR^4 $. Să se arate că dacă $ f $ verifică relația:
\[
  f^2 - 2f + \dr{Id}_{\RR^4} = 0_{\RR^4},
\]
atunci $ f $ este izomorfism și să se determine inversa ei. Să se determine
$ \dr{Ker}f $ și $ \dr{Im}f $ (s-a notat $ f^2 = f \circ f $).

\vspace{0.5cm}

3. (a) Să se diagonalizeze matricea
$
A = \begin{pmatrix}
  0 & 2 & 2 \\
  2 & 0 & 2 \\
  2 & 2 & 0
\end{pmatrix}
$
și apoi să se calculeze $ A^{2018} $.

(b) Fie $ A \in M_6(\RR) $ o matrice nediagonală, cu $ A^2 \neq \alpha I_6 $, pentru
orice $ \alpha \in \RR $, cu proprietatea că:
\[
  A^2 - 5 A + 6I_6 = 0_6.
\]

Să se determine valorile proprii ale matricei $ A $.

\vspace{0.5cm}
\begin{center}
  --------- $\ast\ast\ast\ast\ast$ ---------
\end{center}
\vspace{0.5cm}

\textbf{Seria A -- Numărul 2}

1. Fie subspațiile vectoriale:
\begin{align*}
  V_1 &= \{ (x, y, z) \in \RR^3 \mid 3x - 2y + z = 0 \} \\
  V_2 &= \dr{Sp}\{(1, 1, 2), (0, 1, 2) \}
\end{align*}
ale spațiului vectorial $ \RR^3 $.

Să se determine cîte o bază și dimensiunea pentru $ \RR $-spațiile vectoriale
$ V_1, V_2, V_1 + V_2, V_1 \cap V_2 $. Să se determine coordonatele
vectorului $ v = (15, 8, 7) $ după o bază a lui $ V_1 + V_2 $.

(b) Fie $ V $ un $ \RR $-spațiu vectorial de dimensiune 15. Să se arate
că $ V \simeq \RR^{15} $.

\vspace{0.5cm}

2. (a) Fie $ f : M_2(\RR) \to \RR^3 $ o aplicație liniară astfel încît:
\[
  f \begin{pmatrix} a & b \\ c & d \end{pmatrix} = (a - 2c + d, 0, b - c).
\]

Să se determine cîte o bază și dimensiunea pentru $ \dr{Ker} f, \dr{Im}f $,
matricea lui $ f $ în bazele canonice și să se determine $ V \seq \RR^3 $
astfel încît $ \dr{Im}f \oplus V \simeq \RR^3 $.

(b) Se consideră spațiul vectorial $ \RR^3 $ și o aplicație liniară
$ f : \RR^3 \to \RR^3 $. Să se arate că dacă $ f $ verifică relația:
\[
  f^2 - 3f - \dr{Id}_{\RR^3} = 0_{\RR^3},
\]
atunci $ f $ este izomorfism și să se determine inversa ei. Să se determine
$ \dr{Ker}f $ și $ \dr{Im}f $. Am notat cu $ f^2 = f \circ f $.

\vspace{0.5cm}

3.(a) Să se diagonalizeze matricea:
\[
  A = \begin{pmatrix}
    0 & 3 & 3 \\
    3 & 0 & 3 \\
    3 & 3 & 0
  \end{pmatrix}
\]
și apoi să se calculeze $ A^{2018} $.

(b) Fie $ f : \RR^6 \to \RR^6 $ o aplicație liniară cu proprietatea:
\[
  f^2 - 4f + 3 \dr{Id}_{\RR^6} = 0_{\RR^6}
\]
și $ f^2 \neq \alpha \dr{Id}_{\RR^6} $, pentru orice $ \alpha \in \RR $.

Să se determine valorile proprii ale matricei aplicației liniare.

\newpage

\textbf{Seria D -- Numărul 1}

1. Folosind metoda Gauss-Jordan, să se rezolve sistemul (discuție după $ m $):
\[
  \begin{cases}
    x + y + z &= 1 \\
    x + my + z &= 1 \\
    x + y + mz &= 1
  \end{cases}.
\]

\vspace{0.5cm}

2. (a) Fie $ B = \{ 1, 2 + X \} $ o bază în $ \RR_1[X] $ și $ p = 5 + 3X $.

Aflați coordonatele lui $ p $ în baza $ B $.

(b) Fie subspațiile:
\begin{align*}
  U &= \{ (x, y, z) \in \RR^3 \mid x + y - 2z = 0 \} \\
  V &= \dr{Sp} \{(0, 1, 1), (1, 2, 2) \}.
\end{align*}
Determinați baze și dimensiuni pentru subspațiile $ U, V, U + V $. Decideți
dacă $ U + V $ formează sumă directă și justificați.

\vspace{0.5cm}

3. (a) Fie aplicația liniară:
\[
  f: M_2(\RR) \to \RR_1[X], \quad f \begin{pmatrix} a & b \\ c & d \end{pmatrix} = (a - b) X + d.
\]
Să se determine matricea aplicației în bazele canonice, $ \dr{Ker}f, \dr{Im}f $. Decideți
dacă aplicația este injectivă sau surjectivă și justificați.

(b) Fie $ f : \RR^2 \to \RR^2 $ o aplicație liniară, astfel încît $ f(1, 1) = (1, 0) $
și $ f(2, 1) = (2, 3) $.

Aflați matricea aplicației liniare în baza canonică.

\vspace{0.5cm}

4. Fie matricea $ A = \begin{pmatrix} 2 & 0 \\ 1 & 1 \end{pmatrix} $. Aflați
vectorii și valorile proprii ale matricei.


\vspace{0.5cm}
\begin{center}
  --------- $\ast\ast\ast\ast\ast$ ---------
\end{center}
\vspace{0.5cm}

\textbf{Seria D -- Numărul 2}

1. Folosind metoda Gauss-Jordan, să se rezolve sistemul:
\[
  \begin{cases}
    -2x + y + z &= 1 \\
    x + my + z &= -2 \\
    x + y - 2z &= 4
  \end{cases}.
\]

\vspace{0.5cm}

2.(a) Fie $ B = \{ X, 1 + 2X \} $ o bază a lui $ \RR_1[X] $ și polinomul $ p = 2 + 3X $.
Aflați coordonatele lui $ p $ în baza $ B $.

(b) Fie subspațiile:
\begin{align*}
  U &= \{ (x, y, z) \in \RR^3 \mid x = 2y = 2z \} \\
  V &= \dr{Sp} \{ (1, 2, 1), (0, 1, 0) \}.
\end{align*}

Determinați baze și dimensiuni pentru subspațiile $ U, V, U + V $. Formează $ U $ și
$ V $ sumă directă? Justificați.

\vspace{0.5cm}

3.(a) Fie aplicația liniară:
\[
  f : \RR_2[X] \to M_2(\RR), \quad f(aX^2 + bX + c) = \begin{pmatrix} a & b - c \\ b & 0 \end{pmatrix}.
\]

Să se determine matricea asociată lui $ f $ în bazele canonice, $ \dr{Ker}f $ și $ \dr{Im}f $.
Este $ f $ injectivă sau surjectivă?

(b) Fie $ f : \RR^2 \to \RR^2 $ o aplicație liniară, astfel încît $ f(1, 1) = (1, 0) $ și
$ f(2, 1) = (2, 3) $.

Aflați expresia analitică a aplicației $ f $ ($ f(a, b) = ? $).

\vspace{0.5cm}

4. Fie matricea $ A = \begin{pmatrix} 2 & 4 \\ 1 & 2 \end{pmatrix} $.

Aflați vectorii și valorile proprii.

\newpage

\textbf{Seria E -- Numărul 1}

1. Fie aplicația liniară:
\[
  f : \RR^3 \to \RR^3, \quad f(x, y, z) = (x + y - z, x - y + z, 5x - y + z).
\]
(a) Să se determine o bază și dimensiunea pentru $ \dr{Ker}f, \dr{Im}f $ și
$ \dr{Ker}f + \dr{Im}f $.

(b) Formează $ \dr{Ker}f $ și $ \dr{Im}f $ sumă directă? Justificați.

\vspace{0.5cm}

2. Folosind algoritmul Gram-Schmidt, să se ortogonalizeze vectorii:
\[
  v_1 = (1, 0, 1), v_2 = (-1, 2, 0), v_3 = (0, 1, -1).
\]

\vspace{0.5cm}

3. Fie aplicația liniară:
\[
  f : \RR^2 \to \RR^2, \quad f(x, y) = (x - 2y, 2x - 4y).
\]

Să se găsească valorile proprii ale lui $ f $ și cîte o bază pentru fiecare
subspațiu propriu.

\vspace{0.5cm}

4. Diagonalizați matricea:
$ A = \begin{pmatrix}
  7 & 2 & -2 \\
  2 & 4 & -1 \\
  -2 & -1 & 4
\end{pmatrix} $.

\vspace{0.5cm}
\begin{center}
  --------- $\ast\ast\ast\ast\ast$ ---------
\end{center}
\vspace{0.5cm}

\textbf{Seria E -- Numărul 2}


1. Fie aplicația liniară:
\[
  f : \RR^3 \to \RR^3, \quad f(x, y, z) = (x, y - x, 2x + 3y).
\]
(a) Găsiți o bază și dimensiunea pentru $ \dr{Ker}f, \dr{Im}f $ și $ \dr{Ker}f + \dr{Im}f $.

(b) Este suma $ \dr{Ker}f + \dr{Im}f $ directă? Justificați.

\vspace{0.5cm}

2. Folosind algoritmul Gram-Schmidt, să se ortogonalizeze vectorii:
\[
  v_1 = (0, 1, 1), v_2 = (1, 2, 2), v_3 = (1, 0, 1).
\]

\vspace{0.5cm}

3. Fie aplicația liniară:
\[
  f : \RR^2 \to \RR^2, \quad f(x, y) = (-3x + 10y, -2x + 6y).
\]

Să se găsească valorile proprii ale lui $ f $ și cîte o bază în fiecare
subspațiu propriu.

\vspace{0.5cm}

4. Să se diagonalizeze matricea:
$ A = \begin{pmatrix}
  11 & -5 & 5 \\
  -5 & 3 & -3 \\
  5 & -3 & 3
\end{pmatrix} $.



%%% Local Variables:
%%% mode: latex
%%% TeX-master: "../ag-19-20"
%%% End:
